\chapter{Detalles de implementación e incidencias técnicas}
\label{anx:implementacion}

Este anexo recoge el detalle de implementación que el Capítulo~\ref{cap:desarrollo} resume, y las
incidencias técnicas encontradas durante el desarrollo. Se documentan porque varias de ellas no
producían ningún error visible y habrían invalidado los resultados sin que nada lo delatara.

\section{Entorno de desarrollo}
\label{anx:entorno}

\begin{table}[H]
\centering
\begin{tabular}{|l|l|l|}
\hline
\textbf{Perfil} & \textbf{Servicios} & \textbf{Uso} \\
\hline
\texttt{core} & Bus de eventos y consola & Siempre activo \\
\texttt{stream} & Motor de procesamiento & Desarrollo de detección \\
\texttt{store} & Almacén analítico & Consulta y evaluación \\
\texttt{graph} & Base de datos de grafos & Trabajo sobre el Pilar B \\
\texttt{cti} & Plataforma de inteligencia & Ingesta de CTI \\
\texttt{observability} & Métricas y visualización & Monitorización \\
\hline
\end{tabular}
\caption{Perfiles de despliegue definidos.}
\label{tab:perfiles}
\end{table}

\subsection{Incidencias del entorno}

La construcción del entorno presentó una cadena de incompatibilidades derivadas de emplear
versiones muy recientes del sistema operativo y del intérprete de Python, que se documentan
por su valor práctico:

\begin{itemize}
    \item La distribución de Linux empleada es posterior a las versiones para las que existen
    repositorios oficiales de algunos componentes, lo que obligó a fijar explícitamente el
    repositorio de la última versión con soporte a largo plazo disponible.
    \item La versión de Python del sistema (3.14) carece de paquetes precompilados para las
    bibliotecas de integración con el motor de procesamiento, cuyas dependencias no soportan
    aún esa versión. Se resolvió mediante un entorno virtual independiente con Python 3.12,
    manteniendo el entorno principal para el resto del código.
    \item La biblioteca de integración requiere una máquina virtual Java en el anfitrión, no
    solo en los contenedores, lo que no resulta evidente hasta que falla la primera
    invocación.
\end{itemize}

\section{Selección del subconjunto del esquema}

OCSF define un número elevado de clases de evento, de las cuales solo un subconjunto resulta
pertinente para la telemetría considerada. La selección se realizó a partir de las clases
efectivamente presentes en los conjuntos de datos, y se organizó en tres oleadas de
implementación para permitir que el sistema fuese funcional desde la primera.

\begin{table}[H]
\centering
\begin{tabular}{|l|c|l|}
\hline
\textbf{Clase OCSF} & \textbf{\texttt{class\_uid}} & \textbf{Fuentes que absorbe} \\
\hline
File System Activity & 1001 & Creación y escritura de ficheros \\
Module Activity & 1005 & Carga de bibliotecas dinámicas \\
Scheduled Job Activity & 1006 & Tareas programadas y servicios \\
Process Activity & 1007 & Creación, terminación y acceso a procesos \\
Event Log Activity & 1008 & Manipulación del registro de auditoría \\
Script Activity & 1009 & Ejecución de guiones \\
Account Change & 3001 & Alta y modificación de cuentas \\
Authentication & 3002 & Inicios de sesión \\
Authorize Session & 3003 & Asignación de privilegios especiales \\
Group Management & 3006 & Modificación de pertenencia a grupos \\
Network Activity & 4001 & Conexiones de red \\
DNS Activity & 4003 & Resolución de nombres \\
\hline
\end{tabular}
\caption{Subconjunto de clases OCSF adoptado.}
\label{tab:ocsf}
\end{table}

Merece una mención específica el tratamiento de los eventos de \textbf{registro de Windows}.
El esquema base de OCSF no define una clase específica para esta categoría, por lo que se
optó por asignarlos a la clase más próxima conservando la totalidad del detalle original en
el campo \texttt{unmapped}. La decisión se documenta como tal, dado que una futura extensión
del esquema podría cubrir este caso de forma más precisa.


\section{Diseño del esquema}

El diseño adopta una estrategia \textbf{híbrida}. Los campos que participan habitualmente en
filtros y agregaciones se promocionan a columnas dedicadas con tipos específicos, mientras que
el evento OCSF completo se conserva íntegro en una columna adicional. De este modo se obtiene
el rendimiento de una base de datos columnar sobre los accesos frecuentes sin renunciar a la
flexibilidad de disponer del dato completo cuando se requiere un campo no previsto.

Las decisiones de modelado más relevantes son las siguientes:

\begin{itemize}
    \item La \textbf{clave de ordenación} se establece sobre el par (host, tiempo), que
    corresponde al patrón de acceso dominante en la investigación de incidentes: qué ocurrió
    en una máquina concreta durante una ventana temporal determinada.
    \item El \textbf{particionado} se realiza por mes, lo que facilita la gestión y el purgado
    del almacén por rangos temporales.
    \item Las \textbf{señales de detección} precalculadas durante la normalización se
    materializan como columnas booleanas, lo que reduce las consultas de investigación sobre
    comportamiento adversario a filtros triviales.
    \item Se aplican \textbf{codificaciones específicas} por columna: diferencial más
    compresión para las marcas temporales, que presentan valores próximos entre sí, y
    diccionario para los campos de baja cardinalidad como nombres de host o de clase.
\end{itemize}


\section{Observabilidad}
\label{anx:observabilidad}

La recolección de métricas se realiza mediante un sistema de series temporales que consulta
periódicamente los puntos de exposición de cada componente. Se instrumentaron el bus de
eventos, ambos nodos del motor de procesamiento y el almacén analítico, activando en cada
caso el mecanismo de exposición correspondiente.

Sobre esa base se construyeron tres cuadros de mando, aprovisionados de forma
\textbf{declarativa} —definidos como ficheros de configuración versionados, no configurados
manualmente— para dar cumplimiento al requisito de reproducibilidad RNF1:

\begin{description}
    \item[Salud del \textit{pipeline}] Estado del bus de eventos, volumen por \textit{topic},
    recursos consumidos por el motor de procesamiento y actividad del almacén analítico.
    Verifica el correcto funcionamiento de la infraestructura.
    \item[Panel de seguridad] Distribución de eventos por clase OCSF, señales de detección
    activas, cronología del ataque reconstruida y procesos más frecuentes. Constituye la
    vista orientada al analista.
    \item[Cobertura y calidad de datos] Cobertura de normalización, distribución por canal y
    tipo de evento, conjuntos de datos cargados e indicadores de integridad temporal. Respalda
    las métricas reportadas en la evaluación.
\end{description}


\section{Evaluador de reglas Sigma}
\label{anx:evaluador-sigma}

\subsection{Alcance de los modificadores soportados}

El formato Sigma admite un conjunto amplio de modificadores que alteran la comparación de
valores. Implementarlos todos constituiría un proyecto en sí mismo, por lo que se acotó el
alcance a partir de una medición sobre el repositorio completo de reglas.

\begin{table}[H]
\centering
\begin{tabular}{|l|r|r|}
\hline
\textbf{Modificador} & \textbf{Apariciones} & \textbf{Acumulado} \\
\hline
\texttt{contains}    & 3.362 & 45{,}1\,\% \\
\texttt{endswith}    & 2.494 & 78{,}5\,\% \\
\texttt{all}         &   854 & 90{,}0\,\% \\
\texttt{startswith}  &   431 & 95{,}8\,\% \\
\texttt{re}          &   155 & 97{,}9\,\% \\
\hline
Resto (\texttt{windash}, \texttt{cidr}, otros) & 158 & 100\,\% \\
\hline
\end{tabular}
\caption{Frecuencia de uso de los modificadores de Sigma en el repositorio público.}
\label{tab:modificadores}
\end{table}

Cinco modificadores cubren el 97,9\,\% de los usos reales. El evaluador los implementa
todos, junto con la lógica de composición de condiciones (conjunción, disyunción, negación
y cuantificación). Las construcciones no soportadas se detectan explícitamente y la regla
correspondiente se omite, en lugar de producir una evaluación silenciosamente incorrecta.

\subsection{Optimización}

La implementación inicial reconstruía el árbol de condiciones de cada regla en cada
evaluación. Sobre un conjunto de 2.093 reglas, ello suponía repetir más de 385.000 veces un
trabajo que puede realizarse una sola vez por regla.

Se introdujeron dos optimizaciones:

\begin{itemize}
    \item \textbf{Precompilación}: el árbol de condiciones se construye una única vez al
    cargar la regla.
    \item \textbf{Encaminamiento por tipo de evento}: cada regla declara la categoría de
    registro sobre la que aplica, lo que permite descartar de antemano aquellas que no pueden
    coincidir con el tipo del evento evaluado.
\end{itemize}

El efecto combinado se muestra en la Tabla~\ref{tab:optimizacion}. La comprobación esencial
es que el conjunto de detecciones producidas permaneció \textbf{idéntico}: se trata de una
optimización, no de una aproximación.

\begin{table}[H]
\centering
\begin{tabular}{|l|r|r|r|}
\hline
\textbf{Versión} & \textbf{Tiempo} & \textbf{Mejora} & \textbf{Detecciones} \\
\hline
Sin optimizar & 108{,}6 s & --- & 10 \\
Optimizada    & 1{,}49 s & 73$\times$ & 10 \\
\hline
\end{tabular}
\caption{Efecto de las optimizaciones sobre el conjunto de referencia (184 eventos, 2.093 reglas).}
\label{tab:optimizacion}
\end{table}


\section{Detector de anomalías: protocolo y controles}
\label{anx:anomalias}

\subsection{Protocolo experimental}

El entrenamiento emplea Isolation Forest \cite{isolationforest} sobre la implementación de
scikit-learn \cite{scikitlearn}, y se comparan dos estrategias:

\begin{enumerate}
    \item \textbf{Entrenamiento con ventanas limpias}: únicamente aquellas sin indicios
    objetivos de actividad maliciosa. Corresponde al planteamiento canónico de aprender el
    comportamiento normal.
    \item \textbf{Entrenamiento con contaminación}: todas las ventanas, declarando la
    proporción esperada de anomalías.
\end{enumerate}

Se incorpora asimismo una variante \textbf{supervisada} mediante bosques aleatorios, con y sin
compensación del desbalanceo por SMOTE \cite{smote}.

El conjunto de evaluación se mantiene \textbf{estrictamente separado} del de entrenamiento: se
reserva un escenario completo, para el que se dispone de las alertas de la capa de reglas como
referencia externa. Todos los experimentos se registran en MLflow \cite{mlflow} con sus
parámetros, métricas y semillas, garantizando su reproducibilidad.

\subsection{Controles experimentales}

Junto a los modelos se incorporan dos controles cuya función es delimitar el mérito real de
los resultados:

\begin{itemize}
    \item \textbf{Control de volumen}: un modelo entrenado con una \textbf{única}
    característica, el número de eventos de la ventana. Si alcanzase un rendimiento comparable
    al del modelo completo, quedaría demostrado que la aportación del enfoque multivariante es
    marginal.
    \item \textbf{Ablación de características potencialmente circulares}: se excluyen aquellas
    que participan en la definición de «ventana limpia», para descartar que el modelo
    aproveche información derivada de su propio criterio de entrenamiento.
\end{itemize}

\section{Contaminación del conjunto del laboratorio}
\label{anx:laboratorio}

La revisión del primer conjunto destapó tres \textbf{fuentes de contaminación} que habrían
invalidado el etiquetado sin producir error alguno:

\begin{itemize}
    \item El registro de módulos de PowerShell suponía el 59\,\% del conjunto y era casi
    íntegramente ruido interno del propio marco de pruebas, repartido por igual entre ambas
    clases. Se desactivó y se conservó solo el registro de bloques de guion, que captura el
    contenido real.
    \item Las llamadas de contexto internas del ejecutor de pruebas (\texttt{whoami.exe},
    \texttt{hostname.exe}) suponían el 62\,\% de las muestras maliciosas a nivel de proceso.
    Tras filtrarlas, las muestras maliciosas pasaron de 1.052, con una sola ventana que
    concentraba el 78\,\%, a 147 bien repartidas entre técnicas.
    \item Un servicio del sistema escribía telemetría de compatibilidad de aplicaciones
    durante las ventanas de ataque, lo que añadía 16 filas mal etiquetadas.
\end{itemize}

\section{Incidencias técnicas}
\label{anx:incidencias}

\subsection{Incidencia: sincronización temporal entre fuentes}
\label{sec:incidencia_temporal}

Durante la validación se detectó un error de consecuencias potencialmente graves que no
producía fallo alguno. Los dos canales de telemetría considerados exponen campos temporales
distintos: la telemetría extendida incorpora un campo con la marca de tiempo en formato UTC
además del campo estándar, mientras que la auditoría nativa solo dispone de este último. La
implementación inicial priorizaba el campo aparentemente más preciso cuando estaba
disponible.

El resultado fue un desfase de dieciséis horas entre ambos canales que reordenaba
completamente la secuencia reconstruida: la etapa final del ataque —el borrado del registro de
auditoría— aparecía como la primera. El error no generaba ningún rechazo ni advertencia, y
solo se hizo visible al inspeccionar la coherencia lógica de la cadena de eventos.

La corrección consistió en apoyar la normalización en el campo temporal \textbf{común y
coherente} entre fuentes, en lugar del más específico. La lección resulta generalizable: en un
sistema que correlaciona múltiples fuentes, la homogeneidad del criterio temporal prevalece
sobre la precisión nominal de cada fuente considerada por separado.

Tras la corrección, la secuencia reconstruida presenta un entrelazado coherente de ambas
fuentes con resolución de milisegundos, condición necesaria para la detección de secuencias
del Pilar~A.


\subsection{Incidencia: exposición del servicio de almacenamiento}

Se documenta una segunda incidencia por su carácter no evidente. El servicio de
almacenamiento, ejecutado en contenedor con el puerto correspondiente publicado, rechazaba
las conexiones procedentes del anfitrión pese a responder correctamente a las peticiones
originadas dentro del propio contenedor. El síntoma resultaba engañoso, ya que el error
devuelto sugería un problema de red o de redirección de puertos.

El diagnóstico se alcanzó comprobando que el servicio tampoco respondía a peticiones dirigidas
a su propia dirección de red: escuchaba exclusivamente en la interfaz de bucle local interna.
La solución consistió en declarar explícitamente la interfaz de escucha en la configuración
del servidor.


\subsection{Incidencia: bloqueo del avance temporal por particiones inactivas}
\label{sec:incidencia_watermark}

La causa resultó ser una consecuencia no evidente de la decisión de particionado adoptada. Al
emplear el nombre de host como clave y contener el conjunto de datos un único host, la
totalidad de los eventos se concentra en una sola partición, quedando las restantes vacías.

El motor de procesamiento calcula la marca de agua global como el mínimo de las marcas de agua
de todas sus fuentes. Las particiones que no reciben datos nunca emiten marca de agua alguna,
por lo que el mínimo permanece indefinido y \textbf{el tiempo de evento no avanza jamás}. La
consecuencia es que ninguna ventana temporal llega a considerarse cerrada y el motor no emite
resultado alguno, sin producir error que lo delate.

El diagnóstico se alcanzó por descarte, reduciendo el patrón a un único evento —que debería
haber emitido una coincidencia por cada evento de entrada— y comprobando que tampoco producía
salida, lo que descartaba un problema en la expresión del patrón y señalaba al mecanismo
temporal. La solución consiste en configurar un tiempo de inactividad tras el cual una fuente
sin datos se marca como inactiva y deja de bloquear el cálculo.

Resuelta la incidencia, el patrón de dos eventos emitió correctamente las coincidencias
esperadas, encadenadas y con las marcas temporales originales de la captura. Con ello quedó
validada la vía tecnológica y descartada la necesidad de recurrir a la alternativa prevista
como plan de contingencia.

Esta incidencia ilustra una interacción no trivial entre dos decisiones de diseño
independientes: el particionado por host, adoptado para garantizar el orden de eventos, tiene
como efecto colateral la inutilización del mecanismo de avance temporal cuando el número de
hosts es inferior al de particiones.

\subsection{Incidencia: la hora de los eventos era la de recogida}

El primer día de APT29 no incluye el campo de hora de creación del evento, y la
normalización recurría a la marca del colector, que corresponde al instante de
\textbf{recepción}. El desfase medido sobre los eventos que sí traen la hora real tiene una
mediana de 6,2\,s, un percentil 95 de 74,9\,s y un máximo de 101,2\,s. Lo grave es su
dispersión, de 26,7\,s: los eventos se reordenaban entre sí y no había desplazamiento
constante que corregir.

La corrección establece una cadena de prioridad entre las fuentes de tiempo disponibles, con
la conversión de zona horaria derivada por conjunto de datos, y marca cada evento con la
fuente empleada. Ninguna fuente es fiable en todos los conjuntos: la hora UTC de Sysmon, la
más precisa en APT29, va desviada 19 horas en las capturas de robo de credenciales. Con la
hora corregida desapareció una inversión imposible documentada hasta entonces, en la que el
proceso remoto aparecía tres segundos antes que la conexión que lo provoca: la conexión va
ahora 116\,ms antes. Queda una limitación del dato de origen: el 26,6\,\% de los eventos de
APT29 solo tiene resolución de segundo.


\subsection{Incidencia: la identidad de los eventos y de las alertas}

Durante esta fase se comprobó que el 15\,\% de los eventos normalizados de APT29 era idéntico
byte a byte a otro, porque el evento OCSF no conservaba el número de registro de origen. Dos
ocurrencias distintas con los mismos campos se volvían indistinguibles, y las alertas
duplicadas por un reprocesado no se podían separar de las legítimas. Se resolvió dando
\textbf{identidad determinista} a ambos: el evento conserva su número de registro (o, si la
fuente no lo aporta, el ordinal de la línea en el fichero), la alerta se identifica por el
resumen del evento y la regla, y las tablas del almacén deduplican por esa identidad. La
comprobación final fue recuperar exactamente las 741 alertas documentadas para APT29 y las
cifras de eventos de cada conjunto de datos contadas contra su fichero de origen.

\subsection{Aristas de red del patrón P2}

La construcción de las aristas destapó varios problemas que no producían error alguno.
Bajo la misma clase OCSF de actividad de red conviven las conexiones lógicas de Sysmon y
cada flujo permitido por el cortafuegos de Windows; filtrar por clase daba 3.480 conexiones
WinRM donde el fichero de origen tiene 12. El mapa de direcciones IP a equipos, si cuenta
también los registros entrantes, se corrompe en silencio: un equipo llegó a ganar su propia
dirección por un solo voto. Se resolvió contando solo conexiones salientes, en las que la
dirección de origen es por definición la propia. Y la dirección de bucle local escrita sin
abreviar fabricaba aristas falsas hasta que se normalizó su grafía.

\subsection{Desorden de llegada en modo continuo}

La validación anterior se realizó con fuentes acotadas. En modo continuo, la marca de agua
descarta sin aviso los eventos que llegan con más retraso del admitido, por lo que antes de
lanzar los patrones se \textbf{midió el desorden} de llegada:

\begin{itemize}
    \item Dentro de un mismo conjunto de datos, las alertas llegan con hasta 293\,s de
    retraso respecto a otras posteriores, por efecto de los detectores en paralelo. Con una
    marca de agua de 30\,s se perderían 220 alertas (18\,\%); con 5 minutos, ninguna. La
    cifra está inflada porque el corpus se reproduce decenas de veces más rápido que el
    tiempo real.
    \item Los \textit{topics} intermedios, escritos inicialmente sin ordenar, habrían
    llegado «tarde» en un 74--78\,\%. Se decidió publicarlos en orden de tiempo de evento.
    \item Al reproducir a la vez capturas de distintos años, el desorden entre conjuntos de
    datos alcanza años. Es un artefacto de la reproducción que no existe en producción.
\end{itemize}
